\section{Part III: The Forward Problem - Solving PDEs}
\label{sec:part3}
\lecture{30}{09}{2026}

% Mirrors slides.qmd "Part III" outline. One subsection per outline item.

Parts~\partref{1} and~\partref{2} assembled a complete toolkit on a problem
chosen for its docility. Part~\partref{3} points that toolkit at the
Navier--Stokes equations, where every weakness of the toolkit is provoked at
once: nonlinearity, a global constraint, and complicated geometry.

The order of the argument: what the equations demand
(\cref{subsec:navier-stokes}), how a QTT solver is built
(\cref{subsec:qtt-cfd-frameworks}), why fluids are the right target for this
ansatz rather than an arbitrary one (\cref{subsec:why-fluid-dynamics}), and then
the question everything hinges on --- when does $\chi$ actually stay small
(\cref{subsec:expressivity})?

\subsection{Navier-Stokes Equations}
\label{subsec:navier-stokes}
% Keywords: incompressible/compressible Navier--Stokes, momentum and
% continuity equations, boundary conditions, nondimensionalization,
% Reynolds number, why it is a hard high-dimensional nonlinear PDE.

The incompressible Navier--Stokes equations read
\begin{equation}
  \partial_t \mathbf{v} + (\mathbf{v}\cdot\nabla)\mathbf{v}
    = -\frac{1}{\rho}\nabla p + \nu \nabla^2 \mathbf{v},
  \qquad
  \nabla\cdot\mathbf{v} = 0 .
  \label{eq:navier-stokes}
\end{equation}

They present three difficulties, and --- this is the useful observation --- each
is awkward for a tensor network in a \emph{different} way:

\begin{enumerate}
  \item \textbf{Nonlinear advection}, $(\mathbf{v}\cdot\nabla)\mathbf{v}$:
    products of MPSs, so ranks multiply.
  \item \textbf{A nonlocal pressure constraint}: a global elliptic solve at
    every step.
  \item \textbf{Boundary conditions on complex geometry}: curved interfaces
    couple every length scale at once, hence high rank
    (\cref{subsec:tt-cross-geometry} --- and note that it is the curvature, not
    the sharpness, that costs).
\end{enumerate}

\noindent
\Cref{subsec:tt-cross-geometry} disposed of the third. The first is managed by
rounding after every product. The second is the expensive one.

\paragraph{Why pressure is the hard part.}
\margintext{This is not a tensor-network problem. The globally coupled elliptic
solve dominates runtime in classical incompressible CFD too, which is what makes
the comparison in \cref{subsec:qtt-cfd-results} a fair one.}
\begin{itemize}
  \item Pressure is \textbf{not a dynamical variable}. There is no evolution
    equation for it in~\eqref{eq:navier-stokes}.
  \item It is whatever instantaneously enforces $\nabla\cdot\mathbf{v} = 0$: a
    constraint, not a dynamics.
  \item Taking the divergence of the momentum equation and using
    $\nabla\cdot\mathbf{v} = 0$ to kill the time-derivative and viscous terms
    leaves an elliptic \textbf{Poisson problem},
    \begin{equation}
      \nabla^2 p = -\rho\, \nabla\cdot\!\left[(\mathbf{v}\cdot\nabla)\mathbf{v}\right],
      \label{eq:pressure-poisson}
    \end{equation}%
    whose right-hand side is the only place the nonlinearity enters.
  \item Equation~\eqref{eq:pressure-poisson} couples the whole domain at once:
    the Green's function of $\nabla^2$ has infinite support, so information
    propagates across the mesh within a single step. Every timestep therefore
    contains a globally coupled linear solve. That is where the runtime goes ---
    in classical CFD and in tensor networks alike.
\end{itemize}

\paragraph{Reynolds number and the cascade.}
\begin{itemize}
  \item Nondimensionalize~\eqref{eq:navier-stokes} and exactly one control
    parameter survives: $\mathrm{Re} = UL/\nu$, the ratio of inertial to viscous
    forces.
  \item \textbf{Kolmogorov picture:} energy is injected at the large scale $L$
    and dissipated at the Kolmogorov scale
    $\eta \sim L\,\mathrm{Re}^{-3/4}$, cascading through the scales between.
  \item Resolving both ends in three dimensions therefore needs
    $(L/\eta)^3 \sim \mathrm{Re}^{9/4}$ degrees of freedom --- the estimate
    quoted without justification in \cref{subsec:numerical-methods-intro}.
  \item A commercial aircraft at cruise sits around $\mathrm{Re} \sim 10^7$,
    giving $\sim\!10^{16}$ degrees of freedom. That is the number that makes
    direct numerical simulation impossible.
\end{itemize}

\paragraph{Benchmarks used in these notes.}
\begin{center}
  \begin{tabular}{@{}lp{0.52\textwidth}@{}}
    \toprule
    \textsc{Case} & \textsc{What it tests} \\
    \midrule
    Lid-driven cavity     & closed domain, no immersed geometry; the standard
                            sanity check \\
    Flow past a cylinder  & steady vortex pair, then a K\'arm\'an street;
                            well tabulated \\
    NACA 0040 airfoil     & genuinely non-trivial immersed geometry \\
    Decaying isotropic
    turbulence            & rank behaviour at high $\mathrm{Re}$; the hard
                            case \\
    \bottomrule
  \end{tabular}
\end{center}

\noindent
The compressible cousin of~\eqref{eq:navier-stokes} --- the Euler equations,
with shocks and a nonlinear flux --- is targeted
separately~\cite{Peddinti2025}.

\subsection{QTT Frameworks for CFD}
\label{subsec:qtt-cfd-frameworks}
% Keywords: QTT discretization of velocity/pressure fields, nonlinear
% advection term in TT format, projection/pressure-correction methods,
% lattice Boltzmann in QTT, existing software/frameworks, benchmark flows
% (lid-driven cavity, Burgers' equation).

\paragraph{The encoding.}
\margintext{Every design decision here is inherited from Part~\partref{1}:
interleaving from \cref{subsec:qtt-functions-operators}, low-rank MPOs from the
same place, rounding from \cref{subsec:mps}.}
\begin{itemize}
  \item Each velocity and pressure component becomes \textbf{one MPS} over
    $n = n_x + n_y \,(+\, n_z)$ quantics bits.
  \item Bits are \textbf{interleaved} across coordinates, so each length scale
    stays local along the chain.
  \item Gradient, divergence and Laplacian are MPOs of rank $2$--$4$, built once
    in closed form and reused throughout~\cite{Kazeev2012}.
  \item Everything downstream is MPO application, Hadamard product, and
    rounding. The solver contains no other operation.
\end{itemize}

\paragraph{Chorin projection.}
The scheme is Chorin's projection method~\cite{Chorin1968} carried out in tensor
form~\cite{Peddinti2024}. Three moves per timestep:

\begin{enumerate}
  \item \textbf{Predictor} --- ignore pressure and advance momentum,
    \begin{equation*}
      \mathbf{v}^{*} = \mathbf{v}^{t}
        + \Delta t \left(-(\mathbf{v}^t\cdot\nabla)\mathbf{v}^t
        + \nu\nabla^2\mathbf{v}^t\right).
    \end{equation*}
    The result is generally not divergence-free.
  \item \textbf{Pressure Poisson} --- the DMRG linear solve of
    \cref{subsec:linear-solvers}, and the expensive step,
    \begin{equation*}
      \nabla^2 p = \frac{\rho}{\Delta t}\, \nabla\cdot\mathbf{v}^{*}.
    \end{equation*}
  \item \textbf{Projection} --- subtract the pressure gradient to land back on
    the divergence-free manifold,
    \begin{equation*}
      \mathbf{v}^{t+1} = \mathbf{v}^{*} - \frac{\Delta t}{\rho}\nabla p .
    \end{equation*}
\end{enumerate}

\begin{remark}
  Steps 2 and 3 are a Helmholtz--Hodge decomposition: any vector field splits
  uniquely into a divergence-free part and a gradient, and the Poisson solve is
  how you find the gradient part in order to discard it.
\end{remark}

\begin{remark}[Which Laplacian goes in step 2]
  \margintext{Costs a day to find, ten seconds to fix, and looks exactly like
  \enquote{the solver is a bit inaccurate} the whole time.}
  The operator in step~2 must be $\nabla\cdot\nabla$ \emph{assembled from the
  same difference operators} used for the gradient in step~3 and the divergence
  in step~2. Substitute the compact three-point Laplacian --- which is the
  obvious thing to reach for, since \cref{subsec:qtt-functions-operators} built
  it and it has a lower MPO rank --- and the projection no longer projects.

  The algebra is one line. With $D$ the discrete gradient,
  $\nabla\cdot\mathbf{v}^{t+1} = D\cdot\mathbf{v}^{*} - \Delta t\,(D\cdot D)p$,
  which vanishes identically only if the operator inverted in step~2 was
  $D\cdot D$. With centred differences $D\cdot D \ne L$, and the residual
  divergence is $O(1)$ rather than $O(\text{solver tolerance})$. Measured on one
  step of our demonstrator: $\norm{\nabla\cdot\mathbf{v}}$ falls from $7.7$ to
  $4\times10^{-3}$ when the operators match, and by a factor of a few when they
  do not.
\end{remark}

\paragraph{Geometry and readout.}
\begin{itemize}
  \item \textbf{Geometry.} The smoothed MPS mask of
    \cref{subsec:tt-cross-geometry} is applied elementwise after the predictor;
    ghost-cell inlet and outlet conditions live inside the MPOs.
  \item \textbf{Self-consistent readout} --- the part that makes the method
    usable rather than merely fast:
    \begin{itemize}
      \item a coarse-grained field costs $O(n\chi^3)$, by contracting the fine
        legs against uniform vectors;
      \item a single grid value costs $O(\chi^3)$, by contracting the chain
        against the binary digits of that index.
    \end{itemize}
  \item The dense $2^n$ vector is \textbf{never reconstructed}, at any point in
    the pipeline.
\end{itemize}

\begin{remark}
  Without a compressed readout the whole advantage evaporates: you would
  compress the computation and then decompress it in order to look at it. This
  is a general lesson about compressed solvers, not a detail of this one.
\end{remark}

\paragraph{Results and scaling.}
\label{subsec:qtt-cfd-results}
Incompressible flow past a cylinder~\cite{Peddinti2024}:

\begin{center}
  \begin{tabular}{ll}
    \toprule
    \textsc{Quantity} & \textsc{Value} \\
    \midrule
    Mesh                        & $2^{30}$ cells \\
    Bond dimension              & $\chi = 30$ \\
    Parameters                  & $24{,}200$ \\
    Compression                 & $\sim\!4.4\times 10^{4}$ \\
    Reynolds range (all cases)  & $1.7$--$2230$ \\
    Observed runtime            & $\sim n\chi^{4.1}$ \\
    \bottomrule
  \end{tabular}
\end{center}

\noindent
Runtime is linear in $n$ and therefore \textbf{poly-logarithmic in the mesh
size} $2^n$. The results were validated against Ansys Fluent on the squared
geometries.

\begin{remark}[What the teaching demonstrator reproduces, and what it does not]
  The companion solver in \texttt{notebooks/qtade\_cfd.py} is a few hundred
  lines of NumPy with homogeneous Dirichlet conditions on every wall, no
  ghost-cell inlet or outlet, and no immersed geometry. It is not the code that
  produced the table above, and it is not validated against anything.

  What it does reproduce is the \emph{cost structure}, which is the part worth
  carrying away. On a $64\times64$ grid over 40 timesteps the pressure solve
  takes $82\%$ of the runtime, against the $80.1\%$ reported for the full
  solver. Scaling from $n=5$ to $n=8$ multiplies the cell count by 64 and the
  runtime by about 20 --- sublinear in cells, and well above the ideal $O(n)$,
  the excess being interpreter overhead and local conjugate-gradient iterations
  rather than anything about tensor networks. Take the shape of that curve from
  the demonstrator and the constant from the paper.
\end{remark}

\begin{remark}[Truncation does not respect the constraint]
  SVD truncation is optimal in the Frobenius norm, and the Frobenius norm knows
  nothing about $\nabla\cdot\mathbf{v}=0$. Every rounding step puts a little
  divergence back and the next projection removes it again; the balance is set
  by the rank cap. Measured after 20 steps of the demonstrator:

  \begin{center}
    \begin{tabular}{@{}lr@{}}
      \toprule
      $\chi_{\max}$ & $\norm{\nabla\cdot\mathbf{v}}$ \\
      \midrule
      8  & $1.2\times10^{1}$ \\
      16 & $1.6\times10^{0}$ \\
      32 & $4.1\times10^{-2}$ \\
      48 & $5.9\times10^{-4}$ \\
      \bottomrule
    \end{tabular}
  \end{center}

  \noindent
  Structure-preserving rounding --- truncation that conserves mass, energy or
  the constraint --- is open (\cref{subsec:part2-open-problems}). Until it
  exists, $\chi_{\max}$ is the knob and the divergence is the diagnostic you
  watch.
\end{remark}

\paragraph{Other frameworks in the field.}
\begin{itemize}
  \item \textbf{Reduced-order models} for wall-bounded flows, at a few percent
    of the DNS variable count~\cite{Kiffner2023}.
  \item \textbf{Turbulence PDFs} computed directly in tensor-network form,
    rather than by sampling a reconstructed field~\cite{Gourianov2025}.
  \item \textbf{Vlasov--Poisson} for kinetic plasmas~\cite{Ye2022}.
  \item \textbf{Lattice Boltzmann} in tensor-train form: a different
    discretization, the same compression idea.
\end{itemize}

The common thread is that the ansatz is portable across discretizations, and
that the global solve is always the hard part.

\subsection{Why Fluid Dynamics?}
\label{subsec:why-fluid-dynamics}
% Keywords: computational cost of high-resolution CFD, turbulence and
% scale separation, industrial/scientific relevance, motivation for
% compressed tensor representations.

Two separate questions, worth keeping apart: why fluids are worth the effort,
and why they happen to suit this particular ansatz.

\paragraph{Fluid dynamics is where the cycles go.}
\begin{itemize}
  \item CFD is a top consumer of HPC worldwide: aerodynamics, weather and
    climate, combustion, biomedical flow.
  \item The bottleneck is resolution, and resolution costs
    $\sim\!\mathrm{Re}^{9/4}$ in three dimensions.
  \item Exascale still cannot simulate a full aircraft at flight Reynolds
    number.
  \item Everything cheaper than DNS --- RANS, LES --- is a \textbf{model}, with
    tuned closures and limited transfer between regimes.
\end{itemize}

\paragraph{Why it suits tensor networks.}
\margintext{Gourianov \emph{et al.} (2022) made this argument first and
quantitatively. It is the observation that turned a plausible analogy into a
research programme.}
\begin{itemize}
  \item Turbulence has precisely the structure the ansatz wants: a
    \textbf{cascade} in which well-separated scales couple
    weakly~\cite{Gourianov2022}.
  \item Quantics provides a scale-by-scale representation for free
    (\cref{subsec:qtt-functions-operators}), so the natural hierarchy of the
    physics \emph{is} the natural hierarchy of the ansatz. Weak inter-scale
    coupling is exactly what a small bond dimension encodes.
  \item It is also an honest stress test: nonlinear, multiscale, geometrically
    complex, with unambiguous quantitative benchmarks.
  \item And it is a bridge to hardware: the same encoding is what a
    fault-tolerant quantum CFD algorithm would consume and return.
\end{itemize}

\subsection{Expressivity}
\label{subsec:expressivity}
% Keywords: what functions/fields admit low QTT rank, smoothness vs.
% rank, turbulent vs. laminar flow compressibility, empirical rank
% studies, limits of expressivity.

Everything so far assumed that $\chi$ stays small. This subsection asks when
that is actually true, and it is the most important one in Part~\partref{3}.

\paragraph{The central question.}
Two sub-questions, with different answers:

\begin{itemize}
  \item How does $\chi$ grow with \textbf{resolution}, at fixed physics?
  \item How does $\chi$ grow with \textbf{Reynolds number}?
\end{itemize}

\noindent
The first is benign --- that is the quantics promise of
\cref{subsec:qtt-functions-operators}. The second is the real question, it has
no theory behind it (\cref{subsec:part3-outlook}), and the answer is empirical.

\paragraph{3D turbulence at $\mathrm{Re}_\lambda = 315$.}
\margintext{Pisoni \emph{et al.} (2026): DNS snapshots on $1024^3 = 2^{30}$
points, encoded as 30-tensor trains. The largest such study to date, and the
source of the numbers below.}
\begin{center}
  \begin{tabular}{@{}llp{0.46\textwidth}@{}}
    \toprule
    $\chi$ & \textsc{Retained} & \textsc{What survives} \\
    \midrule
    $1000$ & $2.2\%$  & the full flat inertial range of $E(k)$ \\
    $100$  & $0.03\%$ & enough to run the solver stably for 9--10 turnover
                        times \\
    \bottomrule
  \end{tabular}
\end{center}

\noindent
That is four orders of magnitude of compression with the energy spectrum
essentially intact~\cite{Pisoni2026}.

\paragraph{What breaks first.}
\begin{itemize}
  \item \textbf{The spectrum is the easy diagnostic.} It still looks fine long
    after the flow does not. $E(k)$ is a second-order two-point statistic, and
    second-order statistics are exactly what an $\ell_2$-optimal truncation
    protects. Reporting only $E(k)$ therefore overstates the quality of a
    compression.
  \item The demanding diagnostics are the \textbf{PDFs of velocity increments}
    and the \textbf{flatness}, i.e.\ the intermittency statistics. These are
    high-order and live in the tails, where truncation error concentrates.
  \item \textbf{Bit ordering decides this.} Stacked ordering over-smooths the
    field at low $\chi$; interleaved ordering preserves the non-Gaussian
    statistics.
  \item But below $\chi \approx 500$ the interleaved encoding injects spurious
    small-scale fluctuations. Both encodings fail --- differently.
\end{itemize}

\begin{remark}[The ordering in which diagnostics recover]
  The claim above is easy to check in two dimensions, and worth checking because
  the ordering is so lopsided. Compressing a synthetic intermittent field
  (a $k^{-5/3}$ spectrum modulated by a lognormal multiplicative cascade, true
  flatness $37.6$):

  \begin{center}
    \begin{tabular}{@{}lrrr@{}}
      \toprule
      $\chi$ & $\ell_2$ error & $E(k)$ error, $k<40$ & flatness \\
      \midrule
      16 & $0.61$ & $37\%$   & $119$ \\
      32 & $0.42$ & $13\%$   & $64$ \\
      64 & $0.21$ & $1.4\%$  & $44$ \\
      \bottomrule
    \end{tabular}
  \end{center}

  \noindent
  At $\chi = 64$ the spectrum is right to $1.4\%$ while the pointwise error is
  still $21\%$ and the flatness is $16\%$ high. The diagnostic that survives
  compression longest is therefore the one least able to tell you that the
  compression was safe. State this to anyone who shows you a compressed
  turbulent field and an $E(k)$ plot.
\end{remark}

\begin{remark}[The honest claim]
  What the evidence supports is that turbulence compresses \emph{well at
  moderate} $\mathrm{Re}$. What it does not support is any claim that turbulence
  is low-rank at arbitrary $\mathrm{Re}$ and at every scale. The $\chi = 100$
  ceiling in the current 3D solver is a memory-bound implementation limit, not a
  physical statement.

  A second, cheaper measurement makes the shape of the dependence visible
  without a DNS: synthesize fields with an inertial range of increasing width
  and compress them to a fixed tolerance. With $k_{\max} = 4, 8, 16, 32, 64$ the
  bond dimension needed at $10^{-2}$ goes $14, 26, 41, 68, 128$. The rank tracks
  the \emph{width of the inertial range}, which is the thing $\mathrm{Re}$
  controls --- not the grid, which is the thing quantics already handles.
\end{remark}

The flip side is genuinely encouraging. Random low-rank tensor trains with
$\chi = 10$, cascaded across scales, \textbf{synthesize} divergence-free
turbulent-like snapshots --- correct spectrum, correct flatness --- at
logarithmic cost~\cite{Pisoni2026}. Whatever the ansatz is missing, it
is not the gross statistical structure of turbulence.

\subsection{Outlook}
\label{subsec:part3-outlook}
% Keywords: scaling to 3D/turbulent regimes, adaptive time stepping,
% coupling with machine learning, open numerical-analysis questions.

\paragraph{Solvers and scale.}
\begin{itemize}
  \item \textbf{Higher $\mathrm{Re}$, full 3D.} Lifting the $\chi = 100$ ceiling
    is an engineering problem before it is a physics one.
  \item \textbf{Better solvers.} TT preconditioning, distributed and GPU sweeps,
    mixed precision. Recall that the Poisson solve is $80\%$ of runtime, so
    nothing else is worth optimizing first.
  \item \textbf{Adaptivity.} Choose $\chi(t)$ and $\Delta t$ from a live error
    estimate rather than a fixed cap; tangent-space integrators supply the
    machinery~\cite{Lubich2015}.
  \item \textbf{Structure-preserving truncation.} Maintain
    $\nabla\cdot\mathbf{v} = 0$ and the conserved quantities under SVD rounding.
\end{itemize}

\paragraph{Physics and theory.}
\margintext{The theory gap is the takeaway for an applied-mathematics audience:
there is at present no sharp rank bound for Navier--Stokes solutions as a
function of $\mathrm{Re}$.}
\begin{itemize}
  \item \textbf{Compressible and multiphysics.} Shocks, reacting flows, MHD ---
    regimes where rank behaviour is essentially unmapped.
  \item \textbf{The theory gap.} No sharp rank bounds for Navier--Stokes
    solutions as a function of $\mathrm{Re}$. Contrast this with
    \cref{subsec:mps}, where the one-dimensional area law~\cite{Hastings2007}
    gives exactly such a bound for quantum ground states. That is the missing
    foundation beneath every empirical result quoted above.
  \item \textbf{Hybrid tensor networks and machine learning.} Learned closures,
    learned preconditioners, tensor-network surrogates --- which is
    Part~\partref{4}.
  \item \textbf{Quantum hand-off.} The same MPS encoding is the natural
    interface to future quantum CFD algorithms.
\end{itemize}

\begin{exercise}[Companion notebooks]
  \texttt{03-walkthrough-forward-problem.ipynb} builds the solver above on a
  $64\times64$ grid; \texttt{03-exercises-forward-problem.ipynb} is the hands-on
  half-hour. Keep $n = 6$ unless a cell says otherwise, or you will spend the
  session watching a progress bar.

  \begin{enumerate}
    \item Instrument the advection term and watch the bond dimension of
      $u\,\partial_x u$ \emph{before} rounding. It is $\chi_u \chi_{\partial u}$,
      exactly. Then compare the advective and conservative forms: identical in
      exact arithmetic, not identical under truncation.
    \item Sweep $\chi_{\max}$ and find the smallest value that keeps
      $\norm{\nabla\cdot\mathbf{v}}$ below $10^{-1}$ for thirty steps --- then
      look at what that costs relative to the dense field. On a $64\times64$
      grid the answer is sobering, and the reason is that compression is an
      asymptotic statement.
    \item Put a smoothed cylinder in the flow. One extra Hadamard product per
      step, and no other change. Watch the sign convention on the level set: $q$
      must be \emph{negative} inside the solid, or you will faithfully simulate a
      flow that exists only inside the cylinder.
    \item Compress an intermittent field at several $\chi$ and compare the
      $\ell_2$ error, the spectrum and the flatness. Convince yourself, with your
      own numbers, which diagnostic you are allowed to trust.
  \end{enumerate}
\end{exercise}
